\chapter{Numerical simulations}

\section{The Hamiltonian}

Hamiltonian of the full $\{\ket{0},\ket{1},\ket{e_\mu},\ket{r},\ket{r'}\}$ system takes
\begin{align}
&\ket 0=\ket{5S_{1/2},F=1,m_F=0},\qquad\ket1=\ket{5S_{1/2},F=2,m_F=0},\\
&\ket{e_{F'}}=\ket{6P_{3/2},F',m_F=-1},\qquad F'=1,2,3,\\
&\ket r=\ket{70S_{1/2},m_J=-1/2}, \qquad\ket{r'}=\ket{70S_{1/2},m_J=1/2}.
\end{align}
where the 420-nm beam is assumed to be pure $\sigma_-$ and the 1013-nm beam is assumed to drive the target Rydberg state through the corresponding $\sigma_+$ leg.
The target Hamiltonian before choosing the optical rotating frame is:
\begin{align}
&H_P^{\rm bare}(t)
=-\varepsilon_0\ket0\bra0  + \sum_{F=1}^3 \varepsilon_F\ket{e_{F}}\bra{e_{F}} + (\varepsilon_r-\Delta_R^{\rm bare}(t))\ket r\bra r + (\varepsilon'_r-\Delta_R^{\rm bare}(t))\ket{r'}\bra{r'} \\
& V_{\rm hf}(t) = \frac{\Omega^{\rm bare}_{\rm hf}(t)e^{i\Phi_{\mathrm{hf}}(t)}}{2}\ket1\bra0 + \mathrm{h.c.}, \quad V(t) =
\sum_{F=1}^3 \sum_{a \in \{0,1,r,r'\}}
\left[
\frac{\Omega_{Fa}(t)e^{i\Phi_{Fa}(t)}}{2}\ket{e_{F}}\bra a
+\mathrm{h.c.}\right].
\end{align}
All the laser phase contains the fast rotating terms, i.e.,
\begin{align}
&\Phi_{F0}(t)=\Phi_{F1}(t) = \phi_{\mathrm{420}}(t) - \omega_{\mathrm{420}} t,\\
&\Phi_{Fr}(t) = \Phi_{Fr'}(t) =\phi_{\mathrm{1013}}(t) + \omega_{\mathrm{1013}}t,\\
&\Phi_{\mathrm{hf}}(t) = \varepsilon_0 t + \phi_{\mathrm{hf}}(t).
\end{align}
We first use a unitary transformation $\ket{\psi} \to U\ket{\psi}$, $H = U H U^\dagger + i (\partial_t U) U^\dagger$ to transform the inherent energy level shifts to the driving frequency. A convenient optical-frame transformation matching the convention is
\begin{align}
&U_{\rm opt}(t) =
\exp\left[
-i\sum_{F'=1}^3\theta_{e}(t)\ket{e_{F'}}\bra{e_{F'}}
-i\theta_r(t)\left(\ket r\bra r + \ket{r'}\bra{r'}\right)\right],\\
&i\partial_t U_{\rm opt}(t)U^\dagger_{\rm opt}(t) = \sum_{F'=1}^3\dot\theta_{e}(t)\ket{e_{F'}}\bra{e_{F'}} + \dot\theta_r(t)\left(\ket r\bra r + \ket{r'}\bra{r'}\right).
\end{align}
The phases are chosen so that the 1013-nm couplings are time independent,
while the 420-nm couplings carry one residual programmed phase. In particular,
for the target ladder one may choose the gauge such that
\begin{align}
\theta_{e}(t)-\theta_r(t)= \omega_{\mathrm{1013}}t,
\quad
\theta_{e}(t) =  -\omega_{\mathrm{420}} t,
\end{align}
which solve to the result $\dot\theta_{e}(t) = -\omega_{\mathrm{420}}, \quad \dot\theta_r = (-\omega_{\mathrm{1013}} -\omega_{\mathrm{420}})$.
Defining
\begin{align}
\Delta_e = \varepsilon_{F = 2}  - \omega_{\mathrm{420}}, \quad
\Delta_r = \varepsilon_r - \omega_{\mathrm{1013}} - \omega_{\mathrm{420}}, \quad 
\delta_r = \varepsilon_r' -  \varepsilon_r.
\end{align}
We choose the gauge such that $\Delta_r = 0$.
the Hamiltonian in this rotating frame is
\begin{align}
&H_0(t)
=-\varepsilon_0\ket0\bra0 -\Delta_R^{\rm bare}(t) \left(\ket r\bra r + \ket{r'}\bra{r'}\right) +\delta_r\ket{r'}\bra{r'}
+ \sum_{F=1}^3 (\Delta_e + \eta_F)\ket{e_{F}}\bra{e_{F}}, \label{eq:H_0}\\
& V_{\rm hf}(t) = \frac{\Omega^{\rm bare}_{\rm hf}(t)e^{i(\phi_{\mathrm{hf}}(t) + \varepsilon_0 t)}}{2}\ket1\bra0 + \mathrm{h.c.},\\
& V_{\rm 420}(t) =
\sum_{F=1}^3
\left[
\frac{\Omega_{F1}(t)e^{i\phi_{\mathrm{420}}(t)}}{2}\ket{e_{F}}\bra1
+\mathrm{h.c.}\right] +\sum_{F=1}^3
\left[
\frac{\Omega_{F0}(t)e^{i\phi_{\mathrm{420}}(t)}}{2}\ket{e_{F}}\bra0
+\mathrm{h.c.}\right],\\
& V_{\rm 1013}(t) =\sum_{F=1}^3
\left[
\frac{\Omega_{Fr}(t)e^{i\phi_{\mathrm{1013}}(t)}}{2}\ket{e_{F}}\bra r
+\mathrm{h.c.}\right] + \sum_{F=1}^3\left[
    \frac{\Omega_{Fr'}(t)e^{i\phi_{\mathrm{1013}}(t)}}{2}\ket{e_{F}}\bra{r'}
    +\mathrm{h.c.}\right].
\end{align}

Here $\eta_F$ indicates the hyperfine shift of the $\ket{6P_{3/2},F= 1,2,3}$ state relative to the $\ket{5S_{1/2},F=2,m_F=-1}$ state.
\begin{align}
\eta_{F = 1}=-2\pi\times 51\,{\rm MHz},\qquad
\eta_{F = 2}=0,\qquad
\eta_{F = 3}=2\pi\times 87\,{\rm MHz}.
\end{align}
The $\varepsilon_0$ is the gap between the $\ket{0}$ and $\ket{1}$ states. For $^{87}\mathrm{Rb}$ atoms, $\varepsilon_0 = 2\pi \times  6.835\,{\rm GHz}$.
In the case below, we always assume $V_{\rm hf} = 0$ to cancel the fast rotating terms $e^{i\varepsilon_0 t}$.
After specifying the realistic Hamiltonian, we can proceed to eliminate the large detuning mid-energy levels through rigrous SW-transformation.

\begin{theorem}[Effective 0-1-r-r' Hamiltonian]
\label{thm:stark-elimination}
If the adiabatic and large-detuning conditions
\begin{align}
\left|\frac{\dot\Omega_{Fa}}{\Omega_{Fa}\,\Delta_e}\right|&\ll1,\quad
\left|\frac{\ddot\phi_{\rm 420}}{\Delta_e^2}\right|\ll1,\notag\\
\varepsilon_0,\Delta_e&\gg \eta_F,\Delta_R^{\rm bare},\dot\phi_{\mathrm{420}}(t),
\quad |\delta_r|\ll|\Delta_e|
\end{align}
hold, the second-order block Hamiltonian in the $\ket{0}, \ket{1}, \ket{r}, \ket{r'}$ subspace of the Hamiltonian $H_P^{\rm bare}+V(t)$ is
\begin{align}
H_P^{\rm eff}
&=H_0+
D_0|0\rangle\langle 0|+D_1|1\rangle\langle 1|+D_r|r\rangle\langle r|+D_{r'}|r'\rangle\langle r'|\notag\\
&\quad+
\left(
K_{1r}|1\rangle\langle r|+K_{01}|0\rangle\langle1|+K_{0r}|0\rangle\langle r|
+K_{1r'}|1\rangle\langle r'|+K_{0r'}|0\rangle\langle r'|+K_{rr'}|r\rangle\langle r'|+\mathrm{h.c.}
\right).
\end{align}
with the coefficients given below:
\begin{align}
    D_0 &=-\frac{\sum_{F= 1}^3\Omega^2_{F0}(t)}{4 (\Delta_e + \varepsilon_0)},\quad
    D_1 =-\frac{\sum_{F= 1}^3\Omega^2_{F1}(t)}{4 \Delta_e}, \quad
    D_r =-\frac{\sum_{F= 1}^3\Omega^2_{Fr}(t)}{4  \Delta_e},\quad
    D_{r'} =-\frac{\sum_{F= 1}^3\Omega^2_{Fr'}(t)}{4  \Delta_e},\\
    K_{1r} &=-\frac{\sum_{F= 1}^3\Omega_{Fr}(t)\Omega_{F1}(t)}{4  \Delta_e} e^{-i\phi_{\mathrm{420}}(t)},\\
    K_{01} &=-\left(\sum_{F= 1}^3{\Omega_{F1}(t)\Omega_{F0}(t)} \right) \left(\frac{1}{8  \Delta_e} + \frac{1}{8  \Delta_e + 8 \varepsilon_0} \right),\\
    K_{0r} &=-\left(\sum_{F= 1}^3{\Omega_{F0}(t)\Omega_{Fr}(t)} \right) \left(\frac{1}{8  \Delta_e} + \frac{1}{8  \Delta_e + 8 \varepsilon_0} \right)  e^{-i\phi_{\mathrm{420}}(t)},\\
    K_{1r'} &=-\frac{\sum_{F= 1}^3\Omega_{Fr'}(t)\Omega_{F1}(t)}{4  \Delta_e} e^{-i\phi_{\mathrm{420}}(t)},\\
    K_{0r'} &=-\left(\sum_{F= 1}^3{\Omega_{F0}(t)\Omega_{Fr'}(t)} \right) \left(\frac{1}{8  \Delta_e} + \frac{1}{8  \Delta_e + 8 \varepsilon_0} \right)  e^{-i\phi_{\mathrm{420}}(t)},\\
    K_{rr'} &=-\frac{\sum_{F= 1}^3\Omega_{Fr}(t)\Omega_{Fr'}(t)}{4  \Delta_e}
    \end{align}
\end{theorem}

\begin{proof}
Use the time-dependent Schrieffer-Wolff transformation
\begin{align}
H_{\rm SW}=e^S H e^{-S}+i\dot e^S e^{-S},
\qquad
\ket{\psi_{\rm SW}}=e^S\ket{\psi},
\end{align}
Denote the target Hamiltonian $H_{\varepsilon}=H_0+\varepsilon V$, where $\varepsilon$ is a bookkeeping parameter.
The term $H_0$ and $V= V_{\mathrm 420}+ V_{\mathrm 1013}$ separate with respect to diagonal and off-diagonal subspaces $P\oplus Q$, where $P:= \sum_{a} \ket a\bra a$ and $Q:= \sum_{F'} \ket{e_{F'}}\bra{e_{F'}}$.
The block-off-diagonal part we hope to eliminate is $V_{\mathrm 420}$ and $V_{\mathrm 1013}$.
The SW transformation is given by an anti-Hermitian generator $S(t)=S_1(t)+S_2(t)+\cdots$,
with each $S_n(t) \propto \varepsilon^n$ and block-off-diagonal, i.e.,$P S_n(t)P=Q S_n(t)Q=0$.
Takes the BCH expansion
\begin{align}
H_{\rm SW, \varepsilon}
=
H_0&+ \epsilon\left( V+[S_1,H_0]+i\dot S_1\right)\notag\\
&+\epsilon^2
\left(
[S_2,H_0]+i\dot S_2+[S_1,V]+\frac12[S_1,[S_1,H_0]]+\frac{i}{2}[S_1,\dot S_1]
\right)+O(\epsilon^3).
\end{align}

The $O(\varepsilon)$ off-block part vanish fixes the first order $S_1$ and gives the second-order correction
\begin{align}
V+[S_1,H_0]+i\dot S_1 &= 0 ,\\
P H_{\rm SW, \varepsilon}P
=
H_0 +
\epsilon^2
\left(
[S_1,V]
+
\frac12[S_1,[S_1,H_0]]
+
\frac{i}{2}[S_1,\dot S_1]
\right)
+
O(\epsilon^3) &= H_0 + \frac{\epsilon^2}{2}[S_1,V] + O(\epsilon^3).
\end{align}
Specifically, we take
$
S_1(t)=
\sum_{F=1}^3
\left[
\lambda_{F0}(t)\ket{e_{F}}\bra 0
+\lambda_{F1}(t)\ket{e_{F}}\bra 1
+\lambda_{Fr}(t)\ket{e_{F}}\bra r
+\lambda_{Fr'}(t)\ket{e_{F}}\bra{r'}
\right] - \mathrm{h.c.}.
$
Using $H_0$ in Eq.~(\ref{eq:H_0}), we conclude that:
\begin{align}
&[\,\ket{e_{F}}\bra 0,H_0]=-(\varepsilon_0 +\Delta_e + \eta_F)\ket{e_{F}}\bra 0, \\
&[\,\ket{e_{F}}\bra 1,H_0]=-(\Delta_e + \eta_F)\ket{e_{F}}\bra 1, \\
&[\,\ket{e_{F}}\bra r,H_0]=-(\Delta_e + \eta_F + \Delta_R^{\rm bare})\ket{e_{F}}\bra r,\\
&[\,\ket{e_{F}}\bra{r'},H_0]=-(\Delta_e + \eta_F + \Delta_R^{\rm bare} - \delta_r)\ket{e_{F}}\bra{r'}.
\end{align}
the
$\ket{e_{F}}\bra a$ component is the generator equation
\begin{align}
&\frac{1}{2}\Omega_{F0}(t)e^{i\phi_{\mathrm{420}}(t)}-(\varepsilon_0 +\Delta_e + \eta_F)\lambda_{F0}+i\dot\lambda_{F0}=0, \\
&\frac{1}{2}\Omega_{F1}(t)e^{i\phi_{\mathrm{420}}(t)}-(\Delta_e + \eta_F)\lambda_{F1}+i\dot\lambda_{F1}=0, \\
&\frac{1}{2}\Omega_{Fr} - (\Delta_e + \eta_F + \Delta_R^{\rm bare})\lambda_{Fr}+i\dot\lambda_{Fr}=0,\\
&\frac{1}{2}\Omega_{Fr'} - (\Delta_e + \eta_F + \Delta_R^{\rm bare} - \delta_r)\lambda_{Fr'}+i\dot\lambda_{Fr'}=0.
\end{align}
Writing $\lambda_{Fa}=\tilde\lambda_{Fa}e^{i\phi_{Fa}}$ and $\dot\lambda_{Fa}=(\dot{\widetilde\lambda}_{Fa}+i\dot\phi_{Fa}\,\widetilde\lambda_{Fa})e^{i\phi_{Fa}}$ with
$\left\{\begin{aligned}\phi_{F0} &= \phi_{\mathrm{420}}(t),\\\phi_{F1} &= \phi_{\mathrm{420}}(t),\\ \phi_{Fr} &= \phi_{Fr'} = 0,\end{aligned}\right.$
gives
\begin{align}
\left\{
    \begin{aligned}
&\frac{1}{2}\Omega_{F0}(t)-\left[\varepsilon_0 +\Delta_e + \eta_F + \dot \phi_{\mathrm{420}}(t)\right]\tilde\lambda_{F0}+i\dot{\tilde\lambda}_{F0} = 0,\\
&\frac{1}{2}\Omega_{F1}(t)-\left[\Delta_e + \eta_F + \dot \phi_{\mathrm{420}}(t)\right]\tilde\lambda_{F1}+i\dot{\tilde\lambda}_{F1} = 0,\\
&\frac{1}{2}\Omega_{Fr}(t)-\left[\Delta_e + \eta_F + \Delta_R^{\rm bare} \right]\tilde\lambda_{Fr}+i\dot{\tilde\lambda}_{Fr} = 0,\\
&\frac{1}{2}\Omega_{Fr'}(t)-\left[\Delta_e + \eta_F + \Delta_R^{\rm bare} - \delta_r \right]\tilde\lambda_{Fr'}+i\dot{\tilde\lambda}_{Fr'} = 0.
\end{aligned}
\right.
\end{align}
They are independent ODEs with identical form $\dot{\tilde\lambda}_{Fa}(t)+i\Delta_{Fa}(t)\tilde\lambda_{Fa}(t)=iG_{Fa}(t)$ with
\begin{align}
\left\{
\begin{aligned}
\Delta_{F0}(t) &= \varepsilon_0 +\Delta_e + \eta_F + \dot \phi_{\mathrm{420}}(t),\\
\Delta_{F1}(t) &= \Delta_e + \eta_F + \dot \phi_{\mathrm{420}}(t),\\
\Delta_{Fr}(t) &= \Delta_e + \eta_F + \Delta_R^{\rm bare},\\
\Delta_{Fr'}(t) &= \Delta_e + \eta_F + \Delta_R^{\rm bare} - \delta_r.
\end{aligned}
\right.
\qquad
\left\{
\begin{aligned}
G_{F0}(t) &= \frac{\Omega_{F0}(t)}{2}, \quad G_{F1}(t) = \frac{\Omega_{F1}(t)}{2},\\
G_{Fr}(t) &= \frac{\Omega_{Fr}(t)}{2}, \quad G_{Fr'}(t) = \frac{\Omega_{Fr'}(t)}{2}.
\end{aligned}
\right.
\end{align}
In the case $\varepsilon_0, \Delta_e \gg \eta_F, \Delta_R^{\rm bare}, \dot\phi_{\mathrm{420}}(t), \delta_r$, we can simplify the solution to
\begin{align}
\Delta_{Fr}(t) = \Delta_{Fr'}(t) = \Delta_{F1}(t) = \Delta_e,\quad \Delta_{F0}(t) = \Delta_e + \varepsilon_0, 
\end{align}
For simplicity, we omitted the label and write $\tilde\lambda(t), G(t), \Delta(t)$ uniformly.
Defining $\Theta(t,s)=\int_s^t du\,\Delta(u)$, the solution is $\tilde\lambda(t)
= e^{-i\Theta(t,t_\ast)}\tilde\lambda(t_\ast)+
i\int_{t_\ast}^{t}ds\, e^{-i\Theta(t,s)}G(s)$
for initial condition $\tilde\lambda(t_\ast)$. Define $A(s) :=G(s)/\Delta(s)$ and use the fact that
\begin{align}
&i\int_{t_\ast}^{t}ds\, e^{-i\Theta(t,s)}G(s) = i\int_{t_\ast}^{t}ds\, e^{-i\Theta(t,s)}\Delta(s)A(s)=\int_{t_\ast}^{t}ds\, A(s)\partial_s e^{-i\Theta(t,s)}\notag\\
&=A(s)e^{-i\Theta(t,s)}|_{s = t_\ast}^{t} - \int_{t_\ast}^{t}ds\, e^{-i\Theta(t,s)}\partial_s A(s)= A(t) - A(t_\ast)e^{-i\Theta(t,t_\ast)} -\int_{t_\ast}^{t}ds\, e^{-i\Theta(t,s)} \frac{dA(s)}{ds}.
\end{align}
Plugging back to the solution, we get
\begin{align}
\tilde\lambda(t)
&=
A(t) + e^{-i\Theta(t,t_\ast)}\left(\tilde\lambda(t_\ast) - A(t_\ast)\right) -\int_{t_\ast}^{t}ds\, e^{-i\Theta(t,s)} \frac{dA(s)}{ds}.
\end{align}
We are free to choose the SW transformation with the initial condition $\tilde\lambda(t_\ast) = {G(t_\ast)}/{\Delta(t_\ast)}$ to eliminate the second term.
For the integral term, the fast rotating phase $e^{i \Theta(t,s)}$ eliminates. Rigrously, consider the following identity:
\begin{align}
&I = \int_{t_\ast}^{t}ds\, e^{-i\Theta(t,s)}\frac{dA(s)}{ds}
=\int_{t_\ast}^{t}ds\, \left(\partial_s e^{-i\Theta(t,s)}\right) \frac{1}{i\Delta(s)}\frac{dA(s)}{ds}\\
&=\left.-i \frac{e^{-i\Theta(t,s)}}{\Delta(s)}\frac{dA(s)}{ds}\right|_{t_\ast}^{t} + i \int_{t_\ast}^{t}ds\, e^{-i\Theta(t,s)}\frac{d}{ds}\left[\frac{1}{\Delta(s)}\frac{dA(s)}{ds}\right]\\
\end{align}
Use the condition $\left|{\dot\Omega}/{\Omega\,\Delta}\right|\ll1$ and $\left|{\dot\Delta}/{\Delta^2}\right|\ll1$.\color{red}[Todo: wee need a more rigrous conditions here]\color{black}
\begin{align}
&\lambda_{F0}(t) = \frac{\Omega_{F0}(t)e^{i\phi_{\mathrm{420}}(t)}}{2\Delta_{F0}(t)}, \quad
\lambda_{F1}(t) = \frac{\Omega_{F1}(t)e^{i\phi_{\mathrm{420}}(t)}}{2\Delta_{F1}(t)}, \quad
\lambda_{Fr}(t) = \frac{\Omega_{Fr}(t)}{2\Delta_{Fr}(t)}, \quad
\lambda_{Fr'}(t) = \frac{\Omega_{Fr'}(t)}{2\Delta_{Fr'}(t)}.
\end{align}

Direct multiplication yields
\begin{align}
&\frac{1}{2}P[S_1,V]P  = \sum_{F,F'= 1}^3\sum_{a,b}  \frac{1}{2}P \left[\lambda_{Fa}\ket{e_{F}}\bra{a} -\lambda^*_{Fa}\ket{a}\bra{e_{F}},\frac{\Omega_{F'b}(t)e^{i\phi_{F'b}(t)}}{2}\ket{e_{F'}}\bra{b} + \frac{\Omega_{F'b}(t)e^{-i\phi_{F'b}(t)}}{2}\ket{b}\bra{e_{F'}} \right]P \notag\\
&=- \sum_{F,F'= 1}^3\sum_{a,b}  \frac{\lambda^*_{Fa}\Omega_{F'b}(t)e^{i\phi_{F'b}(t)}}{4}P \left[\ket{a}\bra{e_{F}},\ket{e_{F'}}\bra{b} \right]P + \sum_{F,F'= 1}^3\sum_{a,b}  \frac{\lambda_{Fa}\Omega_{F'b}(t)e^{-i\phi_{F'b}(t)}}{4}P \left[\ket{e_{F}}\bra{a},\ket{b}\bra{e_{F'}} \right]P \notag\\
&=- \sum_{a,b}\sum_{F,F'= 1}^3  \frac{\lambda^*_{Fa}\Omega_{Fb}(t)e^{i\phi_{Fb}(t)}}{4} \ket{a}\bra{b}  - \sum_{a,b} \sum_{F= 1}^3 \frac{\lambda_{Fa}\Omega_{Fb}(t)e^{-i\phi_{Fb}(t)}}{4} \ket{b}\bra{a}.
\end{align}
The $a = b$ terms are the AC Stark shifts, which evaluates to
\begin{align}
D_0 &=  -\left(\sum_{F= 1}^3 \frac{\lambda^*_{F0}\Omega_{F0}(t)e^{i\phi_{F0}(t)}}{4} + \frac{\lambda_{F0}\Omega_{F0}(t)e^{-i\phi_{F0}(t)}}{4}\right) \ket{0}\bra{0} = -\sum_{F= 1}^3\frac{\Omega^2_{F0}(t)}{4 \Delta_{F0}(t)} \ket{0}\bra{0},\\
D_1 &=  -\left(\sum_{F= 1}^3 \frac{\lambda^*_{F1}\Omega_{F1}(t)e^{i\phi_{F1}(t)}}{4} + \frac{\lambda_{F1}\Omega_{F1}(t)e^{-i\phi_{F1}(t)}}{4}\right) \ket{1}\bra{1} = -\sum_{F= 1}^3\frac{\Omega^2_{F1}(t)}{4 \Delta_{F1}(t)} \ket{1}\bra{1},\\
D_r &=  -\left(\sum_{F= 1}^3 \frac{\lambda^*_{Fr}\Omega_{Fr}(t)e^{i\phi_{Fr}(t)}}{4} + \frac{\lambda_{Fr}\Omega_{Fr}(t)e^{-i\phi_{Fr}(t)}}{4}\right) \ket{r}\bra{r} = -\sum_{F= 1}^3\frac{\Omega^2_{Fr}(t)}{4 \Delta_{Fr}(t)} \ket{r}\bra{r},\\
D_{r'} &=  -\left(\sum_{F= 1}^3 \frac{\lambda^*_{Fr'}\Omega_{Fr'}(t)e^{i\phi_{Fr'}(t)}}{4} + \frac{\lambda_{Fr'}\Omega_{Fr'}(t)e^{-i\phi_{Fr'}(t)}}{4}\right) \ket{r'}\bra{r'} = -\sum_{F= 1}^3\frac{\Omega^2_{Fr'}(t)}{4 \Delta_{Fr'}(t)} \ket{r'}\bra{r'}.
\end{align}
The $a \neq b$ terms are equivalent to the effective two-photon Rabi coupling:
\begin{align}
V_{1r} &=  -\frac{1}{2}\sum_{F= 1}^3\left(\frac{\lambda^*_{F1}\Omega_{Fr}(t)e^{i\phi_{Fr}(t)}}{2} + \frac{\lambda_{Fr}\Omega_{F1}(t)e^{-i\phi_{F1}(t)}}{2}\right) \ket{1}\bra{r} + \mathrm{h.c.} \\
&= -\frac{e^{-i\phi_{\rm 420}(t)}}{2}\sum_{F= 1}^3\left(\frac{\Omega_{F1}(t)\Omega_{Fr}(t)}{4\Delta_{F1}(t)}  + \frac{\Omega_{Fr}(t)\Omega_{F1}(t)}{4\Delta_{Fr}(t)} \right)\ket{1}\bra{r} + \mathrm{h.c.} \\
&= -\frac{\sum_{F= 1}^3\Omega_{Fr}(t)\Omega_{F1}(t)}{4  \Delta_e} e^{-i\phi_{\mathrm{420}}(t)} \ket{1}\bra{r} + \mathrm{h.c.}, \\
V_{0r} &= -\frac{1}{2}\sum_{F= 1}^3 \left(\frac{\lambda^*_{F0}\Omega_{Fr}(t)e^{i\phi_{Fr}(t)}}{2} + \frac{\lambda_{Fr}\Omega_{F0}(t)e^{-i\phi_{F0}(t)}}{2}\right) \ket{0}\bra{r} + \mathrm{h.c.} \\
&= -\frac{e^{-i\phi_{\rm 420}(t)}}{2}\sum_{F= 1}^3\left(\frac{\Omega_{F0}(t)\Omega_{Fr}(t)}{4\Delta_{F0}(t)}  + \frac{\Omega_{Fr}(t)\Omega_{F0}(t)}{4\Delta_{Fr}(t)} \right)\ket{0}\bra{r} + \mathrm{h.c.}, \\
&= -\sum_{F= 1}^3{\Omega_{Fr}(t)\Omega_{F0}(t)} \left(\frac{1}{8  \Delta_e} + \frac{1}{8  \Delta_e + 8 \varepsilon_0} \right) e^{-i\phi_{\mathrm{420}}(t)} \ket{0}\bra{r} + \mathrm{h.c.}, \\
V_{01} &=-\frac{1}{2}\sum_{F= 1}^3 \left(\frac{\lambda^*_{F0}\Omega_{F1}(t)e^{i\phi_{F1}(t)}}{2} + \frac{\lambda_{F1}\Omega_{F0}(t)e^{-i\phi_{F0}(t)}}{2}\right) \ket{0}\bra{1} + \mathrm{h.c.} \\
&= -\frac{1}{2}\sum_{F= 1}^3\left(\frac{\Omega_{F0}(t)\Omega_{F1}(t)}{4\Delta_{F0}(t)}  + \frac{\Omega_{F1}(t)\Omega_{F0}(t)}{4\Delta_{F1}(t)} \right)\ket{0}\bra{1} + \mathrm{h.c.}\\
&=-\sum_{F= 1}^3{\Omega_{F1}(t)\Omega_{F0}(t)} \left(\frac{1}{8  \Delta_e} + \frac{1}{8  \Delta_e + 8 \varepsilon_0} \right) \ket{0}\bra{1} + \mathrm{h.c.}
\end{align}
The additional couplings involving $\ket{r'}$ are obtained from the same expression by replacing one of the low-energy labels with $r'$:
\begin{align}
V_{1r'} &=  -\frac{1}{2}\sum_{F= 1}^3\left(\frac{\lambda^*_{F1}\Omega_{Fr'}(t)e^{i\phi_{Fr'}(t)}}{2} + \frac{\lambda_{Fr'}\Omega_{F1}(t)e^{-i\phi_{F1}(t)}}{2}\right) \ket{1}\bra{r'} + \mathrm{h.c.} \notag\\
&= -\frac{\sum_{F= 1}^3\Omega_{Fr'}(t)\Omega_{F1}(t)}{4  \Delta_e} e^{-i\phi_{\mathrm{420}}(t)} \ket{1}\bra{r'} + \mathrm{h.c.},\\
V_{0r'} &= -\frac{1}{2}\sum_{F= 1}^3 \left(\frac{\lambda^*_{F0}\Omega_{Fr'}(t)e^{i\phi_{Fr'}(t)}}{2} + \frac{\lambda_{Fr'}\Omega_{F0}(t)e^{-i\phi_{F0}(t)}}{2}\right) \ket{0}\bra{r'} + \mathrm{h.c.} \notag\\
&= -\sum_{F= 1}^3{\Omega_{Fr'}(t)\Omega_{F0}(t)} \left(\frac{1}{8  \Delta_e} + \frac{1}{8  \Delta_e + 8 \varepsilon_0} \right) e^{-i\phi_{\mathrm{420}}(t)} \ket{0}\bra{r'} + \mathrm{h.c.},\\
V_{rr'} &=  -\frac{1}{2}\sum_{F= 1}^3\left(\frac{\lambda^*_{Fr}\Omega_{Fr'}(t)e^{i\phi_{Fr'}(t)}}{2} + \frac{\lambda_{Fr'}\Omega_{Fr}(t)e^{-i\phi_{Fr}(t)}}{2}\right) \ket{r}\bra{r'} + \mathrm{h.c.} \notag\\
&= -\frac{\sum_{F= 1}^3\Omega_{Fr'}(t)\Omega_{Fr}(t)}{4  \Delta_e} \ket{r}\bra{r'} + \mathrm{h.c.}
\end{align}
\end{proof}

\begin{lemma}[Eliminating the off-resonant $\ket{r'}$ state]
\label{lemma:rprime-elimination}
Let the diagonal energies after eliminating the $6P$ manifold be
\begin{align}
E_0=-\varepsilon_0+D_0,\quad
E_1=D_1,\quad
E_r=-\Delta_R^{\rm bare}+D_r,\quad
E_{r'}=\delta_r-\Delta_R^{\rm bare}+D_{r'}.
\end{align}
If $|E_{r'}-E_a|\gg |K_{ar'}|$ for $a\in\{0,1,r\}$ and the pulse envelopes vary slowly compared with $|E_{r'}-E_a|$, then eliminating $\ket{r'}$ to second order gives an effective $\{\ket0,\ket1,\ket r\}$ Hamiltonian with
\begin{align}
\widetilde D_a&=D_a-\frac{|K_{ar'}|^2}{E_{r'}-E_a},\\
\widetilde K_{ab}
&=
K_{ab}
-\frac12 K_{ar'}K_{br'}^*
\left[
\frac{1}{E_{r'}-E_a}
+
\frac{1}{E_{r'}-E_b}
\right],
\qquad a,b\in\{0,1,r\}.
\end{align}
\end{lemma}

\begin{proof}
Write the post-$6P$ Hamiltonian in the block form $P_3 H P_3+Q_{r'} H Q_{r'}+V_{P_3 r'}+V_{r'P_3}$, where $P_3$ projects onto $\{\ket0,\ket1,\ket r\}$ and $Q_{r'}=\ket{r'}\bra{r'}$. A second Schrieffer-Wolff, equivalently Lowdin, partitioning with generator matrix elements $S_{ar'}=K_{ar'}/(E_a-E_{r'})$ removes $V_{P_3 r'}$ to first order. The second-order block $P_3 H P_3+\frac12P_3[S,V_{P_3 r'}+V_{r'P_3}]P_3$ gives the stated diagonal light shifts and symmetrized off-diagonal coupling renormalization.
\end{proof}

\begin{lemma}[Proportion of $\Omega_{Fa}(t)$, $\Omega_{Fr}(t)$, and $\Omega_{Fr'}(t)$]
    \label{lemma:proportion-of-Omega}
   The proportion of $\Omega_{Fa}(t)$, $\Omega_{Fr}(t)$, and $\Omega_{Fr'}(t)$ are given by
   \begin{align}
    \left\{
    \begin{aligned}
    &\Omega_{F = 1, 0}: \Omega_{F = 2, 0}: \Omega_{F = 3, 0} = \left(\sqrt{\frac{3}{10}} + \sqrt{\frac{2}{15}}\right) \Omega_{420}: -\sqrt{\frac{1}{2}}\Omega_{420}: 0\Omega_{420} \\
    &\Omega_{F = 1, 1}: \Omega_{F = 2, 1}: \Omega_{F = 3, 1} = \left(\sqrt{\frac{3}{10}} -\sqrt{\frac{2}{15}}\right) \Omega_{420}: -\sqrt{\frac{1}{2}} \Omega_{420}: \sqrt{\frac{4}{5}} \Omega_{420} \\
    &\Omega_{F = 1, r}: \Omega_{F = 2, r}: \Omega_{F = 3, r} =\sqrt{\frac{3}{10}} \Omega_{1013}: -\sqrt{\frac{1}{2}} \Omega_{1013}: \sqrt{\frac{1}{5}} \Omega_{1013} \\
    &\Omega_{F = 1, r'}: \Omega_{F = 2, r'}: \Omega_{F = 3, r'} =-\sqrt{\frac{2}{15}} \Omega_{1013}: 0\Omega_{1013}: \sqrt{\frac{1}{5}} \Omega_{1013}
    \end{aligned}
    \right.
   \end{align}
\end{lemma}
\begin{proof}
    The 420-nm coupling is evaluated in the uncoupled $\ket{m_J,m_I}$ basis, where the electric dipole operator preserves $m_I$ and changes $m_J$ by $q=-1$. With $I=3/2$ and CG phases,
    \begin{align}
    \ket1
    &=\frac{1}{\sqrt2}
    \left(
    \ket{m_J=\tfrac12,m_I=-\tfrac12}
    +\ket{m_J=-\tfrac12,m_I=\tfrac12}
    \right),\\
    \ket0
    &=\frac{1}{\sqrt2}
    \left(
    -\ket{m_J=\tfrac12,m_I=-\tfrac12}
    +\ket{m_J=-\tfrac12,m_I=\tfrac12}
    \right).
    \end{align}
    Only two components of each $\ket{e_{F'}}=\ket{6P_{3/2},F',m_F=-1}$ are reached from $m_F=0$ by $\sigma_-$:
    \begin{align}
    \begin{array}{c|cc}
    F'&
    \ket{m_J=-\tfrac32,m_I=\tfrac12}&
    \ket{m_J=-\tfrac12,m_I=-\tfrac12}\\
    \hline
    1&\sqrt{\frac{3}{10}}&-\sqrt{\frac{2}{5}}\\
    2&-\sqrt{\frac{1}{2}}&0\\
    3&\sqrt{\frac{1}{5}}&\sqrt{\frac{3}{5}}
    \end{array}
    \end{align}
    The overall phase of each $\ket{e_{F'}}$ is conventional; the phase choice above matches the signs used for the effective Rabi sum.
    Define
    \begin{align}
    \Omega_{420}(t)
    &:=
    \frac{
    \langle 6P_{3/2},m_J=-\tfrac32,m_I=\tfrac12|
    e\mathbf r\cdot\mathbf E^{(1)}_-(t)
    |5S_{1/2},m_J=-\tfrac12,m_I=\tfrac12\rangle
    }{\sqrt2\,\hbar},\\
    \gamma\,\Omega_{420}(t)
    &:=
    \frac{
    \langle 6P_{3/2},m_J=-\tfrac12,m_I=-\tfrac12|
    e\mathbf r\cdot\mathbf E^{(1)}_-(t)
    |5S_{1/2},m_J=\tfrac12,m_I=-\tfrac12\rangle
    }{\sqrt2\,\hbar}.
    \end{align}
    By the Wigner-Eckart theorem,
    $
    \gamma
    =
    \frac{
    \langle J'=\tfrac32,m_J'=-\tfrac12|d_{-1}|J=\tfrac12,m_J=\tfrac12\rangle
    }{
    \langle J'=\tfrac32,m_J'=-\tfrac32|d_{-1}|J=\tfrac12,m_J=-\tfrac12\rangle
    }
    =\sqrt{{1}/{3}},$
    assuming ideal $\sigma_-$ polarization and no magnetic-field mixing.
    The resulting 420-nm Rabi amplitudes are
    \begin{align}
        \left\{
    \begin{aligned}
    \Omega_{F=1,1}
    &=\left(
    \sqrt{\frac{3}{10}}
    -\sqrt{\frac{2}{5}}\gamma\right)\Omega_{420}
    =  \left(
        \sqrt{\frac{3}{10}}
        -\sqrt{\frac{2}{15}}
        \right)\Omega_{420},\\
    \Omega_{F=1,0}&=\left(
    \sqrt{\frac{3}{10}}
    +\sqrt{\frac{2}{5}}\gamma
    \right)\Omega_{420}
    = \left(
        \sqrt{\frac{3}{10}}
        +\sqrt{\frac{2}{15}}
        \right)\Omega_{420},
    \end{aligned}
    \right.
    \\
    \left\{
    \begin{aligned}
    \Omega_{F=2,1}
    &=
    -\sqrt{\frac{1}{2}}\Omega_{420},\\
    \Omega_{F=2,0}
    &=
    -\sqrt{\frac{1}{2}}\Omega_{420},
    \end{aligned}
    \right.
    \qquad
    \left\{
    \begin{aligned}
    \Omega_{F=3,1}&=
    \left(
    \sqrt{\frac{1}{5}}
    +\sqrt{\frac{3}{5}}\gamma
    \right)\Omega_{420} =
        \sqrt{\frac{4}{5}}\Omega_{420},\\
    \Omega_{F=3,0}
    &=
    \left(
    \sqrt{\frac{1}{5}}
    -\sqrt{\frac{3}{5}}\gamma
    \right)\Omega_{420} =0,
    \end{aligned}
    \right.
    \end{align}
    
    For the target Rydberg state, the 1013-nm coupling uses the $\ket{m_J=-3/2,m_I=1/2}$ component of each $\ket{e_{F'}}$. The off-resonant partner $\ket{r'}$ uses the $\ket{m_J=-1/2,m_I=-1/2}$ component with the same single-electron angular prefactor:
    \begin{align}
    \Omega_{F = 1, r}=\sqrt{\frac{3}{10}}\Omega_{1013}, \quad
    \Omega_{F = 2, r}=-\sqrt{\frac{1}{2}}\Omega_{1013}, \quad
    \Omega_{F = 3, r}=\sqrt{\frac{1}{5}}\Omega_{1013},\\
    \Omega_{F = 1, r'}=-\sqrt{\frac{2}{15}}\Omega_{1013}, \quad
    \Omega_{F = 2, r'}=0, \quad
    \Omega_{F = 3, r'}=\sqrt{\frac{1}{5}}\Omega_{1013}.
    \end{align}
\end{proof}

Using Lemma~\ref{lemma:proportion-of-Omega}, the product sum in Theorem~\ref{thm:stark-elimination} can be represented as parameters $\Omega_{420},\Omega_{1013}$ through:
\begin{align}
    \sum_{F=1}^3{|\Omega_{F1}|^2}
    &=
    \left[\left(\sqrt{\frac{3}{10}}-\sqrt{\frac{2}{15}}\right)^2
    +\frac{1}{2}
    +\left(2\sqrt{\frac{1}{5}}\right)^2\right] {\Omega_{420}^2}=\frac{4}{3}{\Omega_{420}^2},\\
    \sum_{F=1}^3{|\Omega_{F0}|^2}
    &=
    \left[\left(\sqrt{\frac{3}{10}}+\sqrt{\frac{2}{15}}\right)^2
    +\frac{1}{2}
    +0\right] {\Omega_{420}^2}=\frac{4}{3}{\Omega_{420}^2},\\
    \sum_{F=1}^3{|\Omega_{Fr}|^2}
    &=\left[\frac{3}{10}+\frac{1}{2}+\frac{1}{5}\right]{\Omega_{1013}^2}={\Omega_{1013}^2},\\
    \sum_{F=1}^3{|\Omega_{Fr'}|^2}
    &=\left[\frac{2}{15}+0+\frac{1}{5}\right]{\Omega_{1013}^2}=\frac{1}{3}{\Omega_{1013}^2},\\
    \sum_{F= 1}^3\Omega_{Fr}\Omega_{F1}
    &= \left[\sqrt{\frac{3}{10}}\left(\sqrt{\frac{3}{10}} - \sqrt{\frac{2}{15}}\right) + \frac{1}{2} + \sqrt{\frac{4}{5}}\sqrt{\frac{1}{5}}\right]\Omega_{420}\Omega_{1013}=\Omega_{420}\Omega_{1013},\\
    \sum_{F= 1}^3\Omega_{F0}\Omega_{F1}
    &= \left[\left(\sqrt{\frac{3}{10}} - \sqrt{\frac{2}{15}}\right)\left(\sqrt{\frac{3}{10}} + \sqrt{\frac{2}{15}}\right) + \frac{1}{2}\right]\Omega_{420}^2=\frac{2}{3}\Omega_{420}^2,\\
    \sum_{F= 1}^3\Omega_{F0}\Omega_{Fr}
    &=\left[\left(\sqrt{\frac{3}{10}} + \sqrt{\frac{2}{15}}\right)\sqrt{\frac{3}{10}} + \frac{1}{2} + 0\right]\Omega_{420}\Omega_{1013}=\Omega_{420}\Omega_{1013},\\
    \sum_{F= 1}^3\Omega_{F1}\Omega_{Fr'}
    &=\left[-\left(\sqrt{\frac{3}{10}}-\sqrt{\frac{2}{15}}\right)\sqrt{\frac{2}{15}} + 0 + \sqrt{\frac{4}{5}}\sqrt{\frac{1}{5}}\right]\Omega_{420}\Omega_{1013}=\frac{1}{3}\Omega_{420}\Omega_{1013},\\
    \sum_{F= 1}^3\Omega_{F0}\Omega_{Fr'}
    &=\left[-\left(\sqrt{\frac{3}{10}}+\sqrt{\frac{2}{15}}\right)\sqrt{\frac{2}{15}} + 0 + 0\right]\Omega_{420}\Omega_{1013}=-\frac{1}{3}\Omega_{420}\Omega_{1013},\\
    \sum_{F= 1}^3\Omega_{Fr}\Omega_{Fr'}
    &=\left[-\sqrt{\frac{3}{10}}\sqrt{\frac{2}{15}} + 0 + \frac{1}{5}\right]\Omega_{1013}^2=0.
\end{align}

Plugging these sums into Theorem~\ref{thm:stark-elimination} gives the bare $\ket{0}, \ket{1}, \ket{r}, \ket{r'}$ coefficients. Eliminating $\ket{r'}$ with Lemma~\ref{lemma:rprime-elimination} gives the effective $\ket{0}, \ket{1}, \ket{r}$ Hamiltonian:
\begin{align}
H_P^{\rm eff} &= H_0+\widetilde D+\widetilde V_{1r}+\widetilde V_{01}+\widetilde V_{0r},\\
H_0
&=-\varepsilon_0\ket0\bra0 -\Delta_R^{\rm bare}(t)\ket r\bra r,\\
\widetilde D
&=\widetilde D_0\ket0\bra0+\widetilde D_1\ket1\bra1+\widetilde D_r\ket r\bra r,\\
\widetilde V_{ab}
&=\widetilde K_{ab}\ket a\bra b+\mathrm{h.c.},
\qquad (a,b)\in\{(1,r),(0,1),(0,r)\}.
\end{align}
The bare coefficients entering the second elimination are
\begin{align}
D_0 &=-\frac{4}{3} \frac{\Omega_{420}(t)^2}{4(\Delta_e + \varepsilon_0)},\quad
D_1=-\frac{4}{3} \frac{\Omega_{420}(t)^2}{4 \Delta_e},\quad
D_r=-\frac{\Omega_{1013}(t)^2}{4 \Delta_e},\quad
D_{r'}=-\frac{\Omega_{1013}(t)^2}{12 \Delta_e},\\
K_{1r} &=-\frac{\Omega_{420}(t)\Omega_{1013}(t)}{4 \Delta_e} e^{-i\phi_{\mathrm{420}}(t)},\\
K_{01} &=-\frac{2}{3}\Omega_{420}(t)^2 \left(\frac{1}{8  \Delta_e} + \frac{1}{8  \Delta_e + 8 \varepsilon_0} \right),\\
K_{0r} &=-\Omega_{420}(t)\Omega_{1013}(t)\left[\frac{1}{8 \Delta_e} + \frac{1}{8 \Delta_e + 8 \varepsilon_0} \right] e^{-i\phi_{\mathrm{420}}(t)},\\
K_{1r'} &=-\frac{\Omega_{420}(t)\Omega_{1013}(t)}{12 \Delta_e} e^{-i\phi_{\mathrm{420}}(t)},\\
K_{0r'} &=
\frac{\Omega_{420}(t)\Omega_{1013}(t)}{3}
\left[
\frac{1}{8\Delta_e}
+
\frac{1}{8(\Delta_e+\varepsilon_0)}
\right]e^{-i\phi_{\mathrm{420}}(t)},\quad
K_{rr'}=0.
\end{align}
With
\begin{align}
E_0=-\varepsilon_0+D_0,\quad
E_1=D_1,\quad
E_r=-\Delta_R^{\rm bare}+D_r,\quad
E_{r'}=\delta_r-\Delta_R^{\rm bare}+D_{r'},
\end{align}
the $r'$-induced corrections are
\begin{align}
\widetilde D_0&=D_0-\frac{|K_{0r'}|^2}{E_{r'}-E_0},\quad
\widetilde D_1=D_1-\frac{|K_{1r'}|^2}{E_{r'}-E_1},\quad
\widetilde D_r=D_r,\\
\widetilde K_{01}
&=
K_{01}
-\frac12K_{0r'}K_{1r'}^*
\left[
\frac{1}{E_{r'}-E_0}
+
\frac{1}{E_{r'}-E_1}
\right],\\
\widetilde K_{0r}&=K_{0r},\quad
\widetilde K_{1r}=K_{1r}.
\end{align}


 
Considering the Rydberg interaction between $\ket{r} = \ket{nS_{1/2}, m_J= \pm 1/2}$.
The effective Hamiltonian contains one more term
\begin{align}
V_{rr}=\sum_{i<j}\frac{C_{ij}}{R_{ij}^6}n_i^r n_j^r
\end{align}
Here $C_{ij} = 2 \pi \times 874 ~\mathrm{GHz} \cdot \mu m^6 $ is the Coulomb interaction between the atoms $i$ and $j$, and $R_{ij}$ is the distance between the atoms $i$ and $j$.

\section{Adiabatic computation}

A critical step in the Time-optimal (TO) Rydberg CPHASE gate protocol is to design a $\Delta_R(t)$ that allows $\{\ket{00},\ket{01},\ket{10},\ket{11}\}$ to be excited to Rydberg state and then return to $\{\ket{00},\ket{01},\ket{10},\ket{11}\}$ with different accumulated phases.

In analogy, we hope to design a time-dependent pulse $\Omega_R(t),\Omega_{\rm hf,i}(t), \Delta_R(t), \Delta_{\rm hf,i}(t)$ that allows $\{\ket{\mathbf{s}}\}_{\mathbf{s} \in \{\ket{0},\ket{1}\}^n}$ to be excited to Rydberg state and then precisely return to the original state $\{e^{i \phi(\mathbf{s})}\ket{\mathbf{s}}\}_{\mathbf{s} \in \{\ket{0},\ket{1}\}^n}$ with different accumulated phases.

\begin{itemize}
\item For a multi-qubit system, adiabatic computation is more applicable than finding a time-optimal pulse.
\item For different configurations $\ket{s}$, only atoms in the $\ket{1}$ state participate in the Rydberg excitation;
\item atoms in $\ket{0}$ are far detuned from the $\ket{0} \leftrightarrow \ket{r}$ transition and can be dropped from the restricted Rydberg dynamics.
\item This leaves a subsystem of $\ket{1},\ket{r}$ atoms with configuration-dependent interaction terms $V^{(s)}_{ij}$ ($2^n$ configurations in total since each atom can be in $\ket{0}$ or $\ket{1}$).
\item Specifically, we fix a configuration $\ket{s}$, set $J_{ij} := V^{(s)}_{ij}$, and restrict to the $\ket{1},\ket{r}$ subspace.
\end{itemize}

For simplicity, we consider the Pauli basis, i.e.,
consider $X_i = \ket{1}\bra{r}+\ket{r}\bra{1}$, and $Z_i = \ket{1}\bra{1} - \ket{r}\bra{r}$ as the Pauli operators, and a time-dependent Hamiltonian with:
\begin{align}
H(t)
&\doteq \frac{1}{2} \sum_{i} \Omega(t)X_i  + \sum_{i<j} J_{ij} \frac{\mathbb{I}_i - Z_i}{2} \frac{\mathbb{I}_j - Z_j}{2}
+ \sum_{i} \frac{\Delta_R(t)-\Delta_{\rm hf,i}(t)}{2} Z_i\\
&\doteq \frac{1}{2} \sum_{i} \Omega(t)X_i  + \frac{1}{4}\sum_{i<j} J_{ij}  Z_i Z_j + \sum_{i} h_i(t) Z_i := H_X(t) + H_Z(t),
\end{align}
where $\doteq$ means equality up to a scalar energy shift and
\begin{align}
h_i(t)=\frac{\Delta_R(t)-\Delta_{\rm hf,i}(t)}{2}-\frac{1}{4}\sum_{j\neq i}J_{ij}.
\end{align}

\begin{itemize}
\item We consider the case $H(0)=H(T)=H_Z(0)$,  with $\Omega(0) = \Omega(T) = 0$.
\item This indicates that the computational beginning state $\ket{s}$ faithfully returns to the original state $\ket{s}$ with a geometric phase $\phi(s)$.
\item The accumulated phase contains both \textbf{dynamical} and \textbf{geometric} contributions; see the detailed analysis in the note on the geometric phase.
\end{itemize}

Then we analyze how the dynamical and geometric phases are related to the configuration $\ket{s}$, which changes $J_{ij}$ and $h_i$, and whether a global phase can be extracted so that only the relative geometric phase is observable across different configurations.

\section{Atom loss analysis}

The decay rate of the state $\ket{r} = \ket{70S_{1/2}}$ is $\Gamma_r = 1/151.55\,(\mu s)^{-1}$, which gives
\begin{align}
p_{r,i} := \Gamma_r\int_0^T n^r_i(t)\, dt.
\end{align}

The scattering rate from the virtual intermediate state is due to the $420$-nm and $1013$-nm lasers coupling through the $6P_{3/2}$ manifold.
For the clock-state encoding, the Clebsch--Gordan decomposition of the $420$-nm coupling gives the same $4/3$ prefactor as in the CZ-gate model, while the $1013$-nm light shift of the Rydberg state has unit prefactor~\cite{evered2023highfidelity}.
The $6P$ lifetime in rubidium is about $120.7\,\mathrm{ns}$~\cite{gomez2004lifetime}; in the high-fidelity gate simulations of Ref.~\cite{evered2023highfidelity}, a comparable intermediate-state lifetime of $110\,\mathrm{ns}$ is used together with explicit scattering branching ratios.

\begin{theorem}[Intermediate-State Scattering]
\label{thm:intermediate-scattering}
Under condition the mid-state detunings satisfy $|\Delta_e|,\ |\Delta_e+\omega_{01}|\gg |\Omega_{420,\mu}^{(1)}|,\ |\Omega_{420,\mu}^{(0)}|,\ |\Omega_{1013,\mu}|,\ \Gamma_e,\ \|H_g\|$,
then the scattering rate from the intermediate state at time $t$ is approximately
\begin{align}
\dot p_{e,i}^{\rm sc}(t)
\simeq
\Gamma_e
\left[
\frac{4}{3}\left(\frac{\Omega_{420}(t)}{2\Delta_e}\right)^2 n_i^1(t)
+
\frac{4}{3}\left(\frac{\Omega_{420}(t)}{2(\Delta_e+\omega_{01})}\right)^2 n_i^0(t)
+
\left(\frac{\Omega_{1013}}{2\Delta_e}\right)^2 n_i^r(t)
\right],
\end{align}
where the factor $4/3$ is the clock-state Clebsch--Gordan sum for the $420$-nm leg (identical for $\ket0$ and $\ket1$, as derived in the Stark-shift section), the factor $1$ is the $1013$-nm sum, and the $\ket0$ leg carries the extra ground hyperfine detuning $\omega_{01}/2\pi\simeq6.835\,$GHz. The $\ket0$ term is the off-resonant scattering of the global $420$-nm beam on the spectator $\ket0$ atoms; it is suppressed by $(\Delta_e/(\Delta_e+\omega_{01}))^2$ and may be dropped when $\Delta_e\gg\omega_{01}$. Integrating this rate over $t\in[0,T]$ gives the approximation below.
\end{theorem}

\begin{proof}[Proof Sketch]
The rigorous way to eliminate a decaying intermediate manifold is to use the effective-operator formalism for open systems~\cite{reiter2012effective,brion2007adiabatic}. For each spontaneous-emission channel $a$, write
\begin{align}
L_{\mu a}=\sqrt{\eta_{\mu a}\Gamma_e}\ket{a}\bra{e_\mu},
\qquad \sum_a\eta_{\mu a}=1.
\end{align}
Open-system adiabatic elimination gives the effective jump operator
\begin{align}
L_{\rm eff}^{\mu a}=L_{\mu a}H_{\rm NH}^{-1}V_+(t)
=
\sqrt{\eta_{\mu a}\Gamma_e}\,
\ket{a}\left[
\frac{\Omega_{420,\mu}^{(1)}(t)}{2(\Delta_\mu-i\Gamma_e/2)}\bra{1}
+
\frac{\Omega_{420,\mu}^{(0)}(t)}{2(\Delta_\mu+\omega_{01}-i\Gamma_e/2)}\bra{0}
+
\frac{\Omega_{1013,\mu}}{2(\Delta_\mu-i\Gamma_e/2)}\bra{r}
\right],
\end{align}
where the $\ket0$ leg inherits the same $\omega_{01}$-shifted denominator as $\delta_0$.
Thus the instantaneous total scattering rate from site $i$ is
\begin{align}
\dot p_{e,i}^{\rm sc}(t)=
\sum_{\mu,a}\operatorname{Tr}\!\left[
\left(L_{\rm eff}^{\mu a}\right)^\dagger
L_{\rm eff}^{\mu a}\rho_i(t)
\right].
\end{align}
Using $\sum_a\eta_{\mu a}=1$, and dropping the $\ket0$-containing coherences $\rho_{01}^{(i)},\rho_{0r}^{(i)}$ (the $\ket0$ atoms are spectators that carry no coherence with the active $\ket1,\ket r$ subspace in the restricted dynamics) under the same secular approximation, yields
\begin{align}
\dot p_{e,i}^{\rm sc}(t)
=
\Gamma_e\sum_\mu
\left[
\frac{
|\Omega_{420,\mu}^{(1)}(t)|^2\rho_{11}^{(i)}(t)
+
|\Omega_{1013,\mu}|^2\rho_{rr}^{(i)}(t)
+
2\,\operatorname{Re}\!\left[
\Omega_{420,\mu}^{(1)}(t)\Omega_{1013,\mu}^*\,\rho_{1r}^{(i)}(t)
\right]
}{4\Delta_\mu^2}
+
\frac{|\Omega_{420,\mu}^{(0)}(t)|^2\rho_{00}^{(i)}(t)}{4(\Delta_\mu+\omega_{01})^2}
\right]
+O\!\left(\frac{\Omega^2\Gamma_e^3}{\Delta^4},\frac{\Omega^3\Gamma_e}{\Delta^3}\right).
\end{align}
In the large-detuning regime, $(\Delta_\mu^2+\Gamma_e^2/4)^{-1}=\Delta_\mu^{-2}+O(\Gamma_e^2/\Delta_\mu^4)$. Finally, setting $\Delta_\mu\simeq\Delta_e$, replacing $\rho_{00}^{(i)},\rho_{11}^{(i)},\rho_{rr}^{(i)}$ by $n_i^0,n_i^1,n_i^r$, dropping the coherence term under the stated secular or population-only approximation, and using the Clebsch--Gordan sums derived in the Stark-shift section,
\begin{align}
\sum_\mu|\Omega_{420,\mu}^{(1)}|^2=\sum_\mu|\Omega_{420,\mu}^{(0)}|^2=\frac{4}{3}\Omega_{420}^2,\qquad
\sum_\mu|\Omega_{1013,\mu}|^2=\Omega_{1013}^2,
\end{align}
gives the stated estimate.
\end{proof}

Thus we estimate the total intermediate-state scattering probability on site $i$ by
\begin{align}
p_{e,i}^{\rm sc}\simeq
\Gamma_e\int_0^T
\left[
\frac{4}{3}\left(\frac{\Omega_{420}(t)}{2\Delta_e}\right)^2 n_i^1(t)
+
\frac{4}{3}\left(\frac{\Omega_{420}(t)}{2(\Delta_e+\omega_{01})}\right)^2 n_i^0(t)
+
\left(\frac{\Omega_{1013}}{2\Delta_e}\right)^2 n_i^r(t)
\right]\,dt,
\end{align}
where the $\ket0$ term may be dropped in the $\Delta_e\gg\omega_{01}$ limit.
Only part of this scattering is direct atom loss.
Using the branching ratio $\eta_{e\to L}$ into loss channels gives
\begin{align}
p_{e,i}^{\rm loss}\simeq \eta_{e\to L}\,p_{e,i}^{\rm sc}.
\end{align}
For a rough estimate one may take $\Gamma_e\simeq 1/(110\,\mathrm{ns})$ and $\eta_{e\to L}\simeq0.614$ from the error model in Ref.~\cite{evered2023highfidelity}.
The far-detuned scaling $p_e\propto(\Omega/2\Delta_e)^2$ is the usual off-resonant scattering estimate for a weakly saturated optical transition~\cite{grimm2000optical}.

\section{Numerical simulation}

To find the best adiabatic pulse, we use numerical simulation.
We simulate on $L\times L$ square lattice with $L = 2,3,4$. For each scale, we simulate the following Hamiltonian:
\begin{align}
H(t) = \sum_{i=0}^{L^2-1}\left[\frac{\Omega(t)}{2}\,\sigma^x_i - \Delta_i(t)\,n^r_i\right] + \sum_{\langle i,j\rangle} V_{ij}\,n^r_i n^r_j
\end{align}
where $\sigma^x_i=|r\rangle\langle 1|_i+|1\rangle\langle r|_i, n^r_i=|r\rangle\langle r|_i$.

The changeable parameters are $\Omega(t), \Delta(t)$, and $V_{ij} = C_6/r_{ij}^6$ where $C_6=2\pi\times874\ \mathrm{GHz\cdot\mu m^6}$.
We choose lattice spacing $r\in [5\,\mu m, 10\,\mu m]$ for simulation.
This induces a Rydberg interaction strength $V\in [2\pi\times18\ \mathrm{MHz}, 2\pi\times0.874\ \mathrm{MHz}]$.

Besides the adiabatic loss, we also consider Rydberg-state decay and intermediate-state scattering in first-order approximation.
Namely, for each site $i$, we calculate the accumulation of $\ket{r}$ state population $n^r_i(t)$ and $\ket{1}$ state population $n^1_i(t)$ during the evolution.
For the lossless effective two-level trajectory, the total population is conserved, i.e., $n^r_i(t) + n^1_i(t) = 1$.
